\chapter{Richiami di Probabilità}

\noindent
Antes de nada es importante notar la traducción de algunos términos clave entre inglés, italiano y español que son clave para entender la bibliografía y los apuntes del curso:
\begin{itemize}
    \item \tbf{Probability density function} = \underline{Distribuzione/Funzione di densità di probabilità} = Función de densidad de probabilidad
    \item \tbf{Probability mass function} = \underline{Distribuzione/Funzione di massa di probabilità} = Función masa de probabilidad
    \item \tbf{Cumulative distribution function} = \underline{Distribuzione/Funzione di ripartizione}: Función de distribución acumulativa
\end{itemize}
como podemos ver muchas veces en italiano se usa indistintamente el término \textit{distribuzione} o \textit{funzione} para referirse a las funciones de probabilidad. De hecho, a veces se usa el término \tbf{legge} (probabilidad) para referirse a la función de probabilidad, ya que en italiano \textit{legge} significa ley, y en este caso se refiere a la ley que rige el comportamiento de la variable aleatoria.
\begin{observacion}
    Es importante destacar que en otros cursos de Probabilidad y Estadística se ha usado $\mathcal{A}$ para denotar la $\sigma-$álgebra, pero en este curso usaremos $\mathcal{F}$ para evitar confusiones con los eventos.
\end{observacion}

\begin{observacion}
    En muchas ocasiones, usaremos el término \textit{función de densidad} para referirnos a la función masa de probabilidad en variables aleatorias discretas, ya que es un término más general que se aplica tanto a variables aleatorias discretas como continuas, por lo que cuando hablemos de función de densidad, nos referiremos a ambas, y se notará muchas veces también en ambos casos como $P(X)$.
\end{observacion}

\noindent
Utilizamos una moneda trucada donde la probabilidad de que salga cara (fenómeno $A$) es $p\in (0,1)$, por lo que observamos el fenómeno y calcularemos $p$. Tenemos $X=\mathbbm{1}_A$ donde 
\begin{equation*}
    \mathbbm{1}_A(x) =
    \begin{cases}
    1, & \text{si } x \in A,\\
    0, & \text{si } x \notin A
    \end{cases}
\end{equation*}
con $P(A)=p$ y espacio muestral $\Omega=\{0,1\}$. Sabemos que una sola prueba no basta, por lo que haremos muchas (lanzamos $n$ veces la moneda) y tenemos un \tbf{stimatore}

\begin{equation*}
    \overline{X_n}=\frac{1}{n}\sum_{i=1}^n \mathbbm{1}_{Ai}
\end{equation*}

\noindent
Recordemos la desigualdad de Chebyshev para mas tarde probar la \tbf{Legge dei numeri grandi (LGN)}.
\begin{teo}[Disuguaglianza di Markov]
    Sea $X$ una variable aleatoria. Si $X\geq 0$ y $\exists E[X]$, entonces:
    \begin{equation*}
        P[X\geq \varepsilon]\leq \dfrac{E[X]}{\varepsilon} \qquad \forall \varepsilon\in \bb{R}^+
    \end{equation*}
\end{teo}
\begin{prop}(Disuguaglianza Chebyshev)
    Si $X$ es una variable aleatoria tal que $\exists E[X^2]$, se tiene
    \begin{equation*}
        P[|X-E[X]|\geq \delta]\leq \frac{\Var[X]}{\delta^2},\qquad \forall \delta>0
    \end{equation*}
\end{prop}
\begin{proof}
    Se basa en aplicar la desigualdad básica a la variable $(X-E[X])^2$. Para $\varepsilon=\delta^2$, se tiene que:
    \begin{equation*}
        P[|X-E[X]|\geq \delta] = P[(X-E[X])^2\geq \delta^2] \leq \frac{E[(X-E[X])^2]}{\delta^2} = \frac{\Var[X]}{\delta^2}
    \end{equation*}
\end{proof}

\begin{teo}[Legge dei numeri grandi (LGN)]
    Sea $(X_n)_{n\geq 1}$ una sucesión de variables aleatorias tal que
    \begin{enumerate}[label=\arabic*)]
        \item $X_i$ independiente a $X_j$ para $i\neq j$
        \item $E[X_i]=\mu \ \forall i\in \N$
        \item $X_i$ integrables con varianza finita $\Var[X_i]=E[(X_i-E[X_i])^2]=\sigma^2\ \forall i\in \N$
    \end{enumerate}
    entonces
    \begin{equation*}
        P\left[\left|\overline{X}_n-\mu\right|\geq \delta\right]\leq \frac{\sigma^2}{n\delta^2},\qquad \forall \delta>0
    \end{equation*}

    \begin{proof}
    Sea $Z=\overline{X_n}=\frac{1}{n}\sum\limits_{i=1}^n X_i$, entonces $E[Z]=\frac{1}{n}\sum\limits_{i=1}^n E[X_i]=\mu$. Por independencia en $(*)$, tenemos que 
    \begin{equation*}
        E[(X_i-E[X_i])(X_j-E[X_j])]\overset{(*)}{=}E[X_i-E[X_i]]E[X_j-E[X_j]]=0, \qquad \forall i\neq j
    \end{equation*}
    Y finalmente desarrollando obtenemos
    \begin{align*}
        &E(|Z-E[Z]|^2)= E\left[\left|\dfrac{1}{n}\sum\limits_{i=1}^n X_i - E\left[\dfrac{1}{n}\sum\limits_{i=1}^n X_i\right]\right|^2\right] =E\left[\left|\dfrac{1}{n}\sum\limits_{i=1}^n X_i - \dfrac{1}{n}\sum\limits_{i=1}^nE[X_i]\right|^2\right]= \\
        &= E\left[\left|\dfrac{1}{n}\sum\limits_{i=1}^n (X_i - E[X_i])\right|^2\right] = E\left[\dfrac{1}{n^2}\sum\limits_{i=1}^n \sum\limits_{j=1}^n (X_i - E[X_i])(X_j - E[X_j])\right] = \\
        &= \dfrac{1}{n^2}\sum\limits_{i=1}^n \sum\limits_{j=1}^n E[(X_i - E[X_i])(X_j - E[X_j])] \overset{(*)}{=}\dfrac{1}{n^2}\sum\limits_{i=1}^n \Var[X_i] = \dfrac{1}{n^2}\cdot n \sigma^2 = \dfrac{\sigma^2}{n}
    \end{align*}
    donde en $(*)$ sabemos que por independencia (demostrado antes) $E[(X_i - E[X_i])(X_j - E[X_j])]=0$ para $i\neq j$, quedando solo la suma de varianzas. Tras aplicar Chebyshev para la variable $Z$, se tiene lo buscado
    \begin{equation*}
        P\left[\left|Z-E[Z]\right|\geq \delta\right] \leq \frac{Var[Z]}{\delta^2}= \frac{\sigma^2}{n\delta^2}
    \end{equation*}
    \end{proof}
\end{teo}

\section{Teoria della misura}
\begin{definicion}[Spazio misurable]
    Un \tbf{espacio medible} es un par ordenado $(\Omega, \mathcal{F})$ que consiste en
    \begin{itemize}
        \item $\Omega$ (Espacio Muestral): conjunto no vacío que contiene todos los posibles resultados de un experimento o fenómeno.
        \item $\mathcal{F}$ ($\sigma$-álgebra): familia (o colección) de subconjuntos de $\Omega$ que satisface tres axiomas específicos que garantizan que las operaciones lógicas entre sucesos sigan siendo sucesos.
    \end{itemize}
\end{definicion}

\begin{definicion}[$\sigma-$algebra]
    Sea $\mathcal{F}$ una colección de subconjuntos de un conjunto no vacío $\Omega$, decimos que es una $\sigma-$álgebra si cumple
    \begin{multicols}{2}
        \begin{enumerate}
            \item $\emptyset \in \mathcal{F}$
            \item Si $A\in \mathcal{F} \implies A^c \in \mathcal{F}$
            \item Si $(A_n)_{n\geq 1} \in \mathcal{F} \implies \bigcup_{n\geq 1} A_n \in \mathcal{F}$
        \end{enumerate}
    \end{multicols}
\end{definicion}

\begin{definicion}[Spazio di probabilità]
    Un \tbf{espacio de probabilidad} es una tripla $(\Omega, \mathcal{F}, P)$ donde
    \begin{enumerate}[label=\arabic*)]
        \item $(\Omega, \mathcal{F})$ es un espacio medible.
        \item $P:\mathcal{F}\to [0,1]$ es una función llamada \tbf{medida de probabilidad} que asigna a cada suceso un número real entre 0 y 1, cumpliendo los Axiomas de Kolmogorov
        \begin{itemize}
            \item Normalización (N): la probabilidad del espacio total es 1, es decir, $P(\Omega) = 1$
            \item $\sigma$-aditividad (A): para cualquier secuencia de sucesos $\{A_i\}_{i \geq 1}$ que sean disjuntos entre sí (es decir, $A_i \cap A_j = \emptyset$ para $i \neq j$), la probabilidad de su unión es la suma de sus probabilidades individuales
            $$P\left(\bigcup_{n\geq 1} A_n\right) = \sum_{n\geq 1} P(A_n)$$
        \end{itemize}
    \end{enumerate}
\end{definicion}

\noindent
\tbf{Quesito}. Perché non posso sempre scegliere $\mathcal{F}=\mathcal{P}(\Omega)$? (l'insieme delle parti di $\Omega$) Esto es posible siempre que $\Omega$ numerable.

\begin{teo}[Vitali 1905]
    The power set is too large. Let $\Omega = \{0,1\}^{\mathbb{N}} $ be the sample space for infinite coin tossing. Then there is no mapping $P: \mathcal{P}(\Omega) \to [0,1]$ with the properties
    \begin{itemize}
        \item (N) Normalisation. $P(\Omega) = 1$.
        \item (A) $\sigma$-Additivity. If $A_1, A_2, \ldots$ are pairwise disjoint, then
        $$P\left(\bigcup_{n\geq 1} A_n\right) = \sum_{n\geq 1} P(A_n).$$
        (The probabilities of countably many incompatible events can be added up.)
        \item (I) Flip invariance. For all $A \subseteq \Omega$ and $n \geq 1$ one has $P(T_n A) = P(A)$; here
        $$T_n: \Omega \to \Omega, \quad \omega = (\omega_1, \omega_2, \ldots) \mapsto (\omega_1, \ldots, \omega_{n-1}, 1 - \omega_n, \omega_{n+1}, \ldots)$$
        is the mapping from $\Omega$ onto itself that inverts the result of the nth toss, and \\ $T_n A = \{T_n(\omega) : \omega \in A\}$ is the image of $A$ under $T_n$. (This expresses the fairness of the coin and the independence of the tosses.)
    \end{itemize}
\end{teo}

\noindent
Es importante destacar que la propiedad de \textit{Flip invariance} simplemente nos indica que la definición de la función de probabilidad $P$ debe tener sentido, ya que puede existir $P$ que cumpla (N) y (A) pero no (I), es decir, que no sea consistente con la naturaleza del experimento aleatorio que estamos modelando. En este caso, lo que se está expresando en concreto es que si en un conjunto $A$ de secuencias de lanzamientos de moneda, invertimos el resultado de un lanzamiento específico (el n-ésimo), la probabilidad asignada a ese conjunto no debería cambiar, ya que hemos cambiado un resultado de manera simétrica y justa en todos los casos.

% //TODO: ejemplos de espacios de probabilidad con diferentes sigma-algebras

\begin{prop}
    Sea $(\Omega, \mathcal{F}, P)$ un espacio de probabilidad, allora
    \begin{enumerate}[label=\arabic*)]
        \item $(A_n)_{n\geq 1} \in \mathcal{F}$ tal que $A_n\subseteq A_{n+1} \ \forall n  \implies P(\bigcup\limits_{n\geq 1} A_n)=\lim\limits_{n\to \infty} P(A_n)$
        \item $(B_n)_{n\geq 1} \in \mathcal{F}$ tal que $B_{n+1}\subseteq B_n \ \forall n  \implies P(\bigcap\limits_{n\geq 1} B_n)=\lim\limits_{n\to \infty} P(B_n)$
    \end{enumerate}
    % //TODO: proof
\end{prop}

\section{Variabili aleatorie}

\begin{observacion}
    Trabajaremos siempre con un espacio de probabilidad $(\Omega, \mathcal{F}, P)$ fijo.
\end{observacion}

\begin{definicion}[Variabile aleatoria]
    Una \tbf{variable aleatoria} $X$ es una función medible sobre un espacio de probabilidad. Es decir, una función:
    $$X:(\Omega, \mathcal{F}, P) \rightarrow (\R, \mathcal{B}(\R)$$
    tal que la imagen inversa de cualquier conjunto de Borel es medible. Esto se caracteriza por:
    $$X^{-1}(B) \in \mathcal{F} \qquad \forall B \in \mathcal{B}(\R)$$
\end{definicion}

\begin{observacion}
    Esta definición de variable aleatoria se puede aplicar a cualquier espacio medible de llegada, no solo a $\R$ con su $\sigma-$álgebra de Borel.
\end{observacion}

\begin{definicion}[Legge di probabilità]
    Sea $X:\Omega \to \mathbb{R}$ una variable aleatoria, la \textbf{ley de probabilidad} de $X$ es la medida de probabilidad $P_X$ definida en $(\mathbb{R}, \mathcal{B}(\mathbb{R}))$ como la siguiente función
    \Func{P_X}{\left({\R}, \mathcal{B}\right)}{[0,1]}{B}{P(X^{{-1}}(B))}
    también se nota como $\mu_X(B)=P_X(B)= P[X\in B]$.
\end{definicion}

\begin{definicion}[Funzione di distribuzione (cumulativa)/funzione di ripartizione]
    Se define la \tbf{función de distribución (acumulativa)} de una variable aleatoria como:
    \Func{F_X}{\R}{[0,1]}{x}{P_X\left(]-\infty, x]\right)}
    también se nota como $F_X(x)=P(X\leq x)=\mu_X(\left]-\infty, x\right])$.
\end{definicion}

\begin{coro}
    Sea $X$ una variable aleatoria, y sea $(z_n)_{n\geq 1}$ una sucesión de números reales tal que $z_n \geq z_{n-1} \geq \dots\geq z_1=z$ y $\lim\limits_{n\to \infty} z_n = z$, entonces
    \begin{equation*}
        \lim_{n\to \infty} F_X(z_n) = F_X(z)
    \end{equation*}
\end{coro}

\begin{teo}[Probability rules]
    Sea $(\Omega, \mathcal{F}, P)$ un espacio de probabilidad, se cumplen las siguientes propiedades para eventos arbitrarios $A, B, A_1, A_2, \dots \in \mathcal{F}$.
    \begin{enumerate}[label=(\alph*)]
        \item $P(\emptyset)=0$
        \item \tbf{Aditividad finita}: $P(A\cup B) + P(A\cap B)=+P(A)+P(B)$, y en particular $P(A)+P(A^c)=1$
        \item \tbf{Monotonía}: si $A\subseteq B$ entonces $P(A)\leq P(B)$
        \item \tbf{$\sigma-$subaditividad}: $P\left(\bigcup\limits_{i=1}^{\infty} A_i\right) \leq \sum\limits_{i=1}^{\infty} P(A_i)$
        \item \tbf{$\sigma-$continuidad}: si $A_n \nearrow A$ o $A_n \searrow A$ (es decir, los $A_n$ son crecientes con unión $A$, o decrecientes con intersección $A$), entonces $P(A_n) \to P(A)$ cuando $n\to \infty$
    \end{enumerate}
\end{teo}

\subsection{Propietà funzione di ripartizione}

\begin{prop}
    Sea $X$ una variable aleatoria, entonces su función de distribución $F_X$ cumple
    \begin{enumerate}[label=\arabic*)]
        \item $F_X$ es monótona no decreciente
        \item $\lim\limits_{x\to -\infty} F_X(x) = 0$ y $\lim\limits_{x\to +\infty} F_X(x) = 1$
        \item si $z_n \searrow z$, entonces $\lim\limits_{n\to \infty} F_X(z_n) = F_X(z)$
        \item se $z_n \nearrow z$, entonces $\lim\limits_{n\to \infty} F_X(z_n) = F_X(z) \Longleftrightarrow P(X=z)=0$
    \end{enumerate}
\end{prop}

\begin{definicion}[Variabile aleatoria discreta]
    Sea una variable aleatoria $X:\Omega\to \R$ con $(\Omega, \mathcal{F},P)$ un espacio de probabilidad, se dice que es \tbf{discreta} si 
    $$ X(\Omega)=\{X(w): w\in \Omega\}$$
    es a lo máximo numerable.
    
\end{definicion}

\begin{definicion}[Variabile aleatoria continua]
    Una variable aleatoria $X$ se dice \tbf{continua} si 
    \begin{equation*}
        P(X=x)=0, \qquad \forall x\in \R
    \end{equation*}
    o equivalentemente, si su función de distribución $F_X$ es continua en todo punto de $\R$.
\end{definicion}


\begin{definicion}[Misura]
    Una \tbf{medida} es cualquier función que asigna un ``tamaño'' o ``peso'' a los conjuntos de una $\sigma$-álgebra. Para que sea una medida, solo debe cumplir tres reglas básicas
    \begin{enumerate}[label=\arabic*)]
        \item \tbf{No negatividad}: el tamaño no puede ser negativo ($\mu(A) \geq 0$).
        \item \tbf{Vacío nulo}: el conjunto vacío mide cero ($\mu(\emptyset) = 0$).
        \item \tbf{$\sigma$-aditividad}: el tamaño de la unión de piezas disjuntas es la suma de los tamaños de las piezas
        \begin{equation*}
            \mu\left(\bigcup_{i=1}^{\infty} A_i\right) = \sum_{i=1}^{\infty} \mu(A_i)
        \end{equation*}
    \end{enumerate}
\end{definicion}

\begin{definicion}[Misura assolutamente continua]
    Sea un espacio medible $(\Omega, \mathcal{F})$ y dos medidas $\mu$ y $\nu$ definidas en él. Decimos que $\mu$ es \tbf{absolutamente continua} respecto a $\nu$ (y lo denotamos como $\mu \ll \nu$) si se cumple que
    \begin{equation*}
        \text{si } \nu(A)=0  \implies \mu(A)=0 \quad \forall A\in \mathcal{F}
    \end{equation*}
\end{definicion}

\begin{teo}[Teorema de Radon-Nikodym]
    Sea $(\Omega, \mathcal{F})$ un espacio medible y $\mu$ y $\nu$ dos medidas definidas en él. Si $\mu$ es absolutamente continua respecto a $\nu$ (es decir, $\mu \ll \nu$), entonces existe una función integrable $f:\Omega \to [0, +\infty]$ ($f\in L_{\mu}^1$) tal que
    \begin{equation*}
        \mu(A) = \int_A f(x) \, d\nu(x)
    \end{equation*}
\end{teo}


\begin{prop}
    Sea $X$ una v.a. ass. continua con función de densidad $f_X$, entonces se tiene
    \begin{equation*}
        E[g(X)]=\int_{\R} g(z)f_X(z)dz
    \end{equation*}
    donde $g:\R \to \R$ positiva.
    \begin{proof}
        \tbf{Paso 1}: Funciones Indicadoras (El átomo). Supongamos primero que $g$ es una función indicadora sobre un conjunto de Borel $A$, es decir
        $$g(x) = \mathbbm{1}_A(x) = \begin{cases} 1 & \text{si } x \in A \\ 0 & \text{si } x \notin A \end{cases}$$
        
        \begin{itemize}
            \item Lado izquierdo: $E[g(X)] = E[\mathbbm{1}_A(X)]$. Por definición, la esperanza de una indicadora es la probabilidad del suceso: $P(X \in A)$.
            \item Lado derecho: $\int_{\mathbb{R}} \mathbbm{1}_A(z) f_X(z) dz = \int_A f_X(z) dz$. Por definición de variable absolutamente continua, esta integral es exactamente $\mu_X(A)$, que es $P(X \in A)$.
        \end{itemize}
        Conclusión: la igualdad se cumple para funciones indicadoras. \\
        \tbf{Paso 2}: Funciones Simples (Linealidad). Una función simple es una combinación lineal de indicadoras $g(x) = \sum_{i=1}^n a_i \mathbbm{1}_{A_i}(x)$.Debido a que tanto la Esperanza ($E$) como la Integral ($\int$) son operadores lineales, podemos sacar las constantes y separar las sumas
        \begin{align*}
            &E\left[ \sum a_i \mathbbm{1}_{A_i}(X) \right] = \sum a_i E[\mathbbm{1}_{A_i}(X)] = \sum a_i \int_{A_i} f_X(z) dz =\\
            &= \int_{\mathbb{R}} \left( \sum a_i \mathbbm{1}_{A_i}(z) \right) f_X(z) dz=\int_{\R}g(z)f_X(z)dz
        \end{align*}
        \tbf{Paso 3}: Funciones No Negativas (Límite Monótono). Aquí es donde ocurre la "magia" del análisis. Si $g$ es una función medible y positiva ($g \geq 0$), existe una sucesión de funciones simples $(g_n)_{n \geq 1}$ que crece hacia $g$ punto a punto ($g_n \nearrow g$).  \\
        \noindent
        Aplicamos el Teorema de la Convergencia Monótona (TCM): como $g_n \nearrow g$, entonces \mbox{$E[g_n(X)] \to E[g(X)]$}. Como ya demostramos en el Paso 2 que la igualdad vale para cada $g_n$ tenemos lo siguiente
        $$E[g_n(X)] = \int_{\mathbb{R}} g_n(z) f_X(z) dz$$
        Al tomar el límite en ambos lados, el TCM nos permite "meter" el límite dentro de la integral
        $$\lim_{n \to \infty} \int_{\mathbb{R}} g_n(z) f_X(z) dz = \int_{\mathbb{R}} g(z) f_X(z) dz$$
    \end{proof}
\end{prop}


\begin{observacion}
    Si $f\in L_{\mu}^1$ se debe cumplir que la integral del valor absoluto de la función sea finita, es decir 
    $$\int_{\Omega} |f(\omega)| \, d\mu(\omega) < \infty$$
    Si esta condición se cumple, decimos que $f$ es integrable.
\end{observacion}

\begin{observacion}
    El tipo de integrales con las que estamos trabajando son integrales de Lebesgue, que generalizan las integrales de Riemann y permiten integrar funciones más generales en espacios medibles.Si la medida $\nu$ tiene una función de densidad $\rho(x)$ (como la densidad de una normal o una binomial), entonces se cumple que
    $$d\nu(x) = \rho(x) \, dx$$
    En este caso, la integral "extraña" se convierte en una integral común
    $$\int_A f(x) \, d\nu(x) = \int_A f(x) \rho(x) \, dx$$
\end{observacion}

\begin{definicion}[Variabile aleatoria assolutamente continua]
    Una variable aleatoria $X$ se dice \tbf{absolutamente continua} si 
    \begin{equation*}
        \mu_X(A)=\int_A f_X(x) \, dx, \qquad \forall A\in \mathcal{B}(\R)
    \end{equation*}
    para una función $f_X:\R \to [0, +\infty]$ positiva e integrable sobre $\R$. Modo alternativo $\mu_x \ll Leb$.
\end{definicion}

\begin{observacion}
    \begin{enumerate}[label=\arabic*)]
        \item $\int\limits_{\R} f(x)dx=\mu_X(\R)=P(X\in \R)=P(\Omega)=1$ donde $f$ es la función de densidad de $X$.
        \item si $X$ es absolutamente continua $\implies$ $X$ es continua. 
    \end{enumerate}
\end{observacion}

\section{Richiami di combinatoria}
Estudiaremos las diferentes formas de contar arreglos, selecciones o agrupaciones de elementos de un conjunto finito. Consideremos un conjunto $\Omega$ con $n$ elementos distintos, es decir, $\Omega=\{1,2,\dots,n\}$, con por ejemplo $n$ bolas enumeradas.
\begin{enumerate}[label=\arabic*)]
    \item \tbf{Disposizioni di classe $k$ di $n$ elementi con ripetizioni con ordine}: las \textit{disposiciones de clase $k$ con repeticiones y con orden} son los vectores $(x_1,x_2,\dots,x_k)$ con $x_i \in \Omega$. Forman el siguiente conjunto 
    \begin{equation*}
        D_n^k=\{(x_1, \dots, x_k): x_i \in \Omega\} \text{ y su cardinal es } |D_n^k|=n^k
    \end{equation*}
    
    \item \tbf{Disposizioni di classe $k$ di $n$ elementi senza ripetizioni con ordine}: las \textit{disposiciones de clase $k$ sin repeticiones y con orden} son los vectores $(x_1,x_2,\dots,x_k)$ con $x_i \in \Omega$ y $x_i \neq x_j$ para $i\neq j$. Forman el siguiente conjunto
    \begin{equation*}
        D_{sr,n}^k=\{(x_1, \dots, x_k): x_i \in \Omega, x_i \neq x_j \text{ per } i\neq j\}, \text{ y su cardinal es } |D_{sr,n}^k|=\dfrac{n!}{(n-k)!}
    \end{equation*} 
    \item \tbf{Disposizioni di classe $k$ di $n$ elementi senza ripetizioni senza ordine (=combinazioni)}: las \textit{combinaciones de clase $k$ sin repeticiones y sin orden} son las anteriores pero sin importar el orden, es decir, los subconjuntos $\{x_1,x_2,\dots,x_k\} \sim (y_1,y_2, \dots, y_k)$ se diferencian en una permutación, por lo que forman el cardinal es 
    \begin{equation*}
        C_n^k=\dfrac{D_{sr,n}^k}{k!}=\dfrac{n!}{k!(n-k)!}=\binom{n}{k}
    \end{equation*}
    
\end{enumerate}

\section{Distribuciones Discretas}

\noindent
En esta sección, recordaremos algunas de las distribuciones de probabilidad discretas más importantes que se utilizan en estadística y probabilidad.

\subsection{Distribución uniforme discreta}

\begin{definicion}
    Sea $X$ una variable aleatoria. Se dice que la variable aleatoria $X$ se distribuye uniformemente alrededor de los puntos $x_1,\dots, x_n$ si dichos valores son equiprobables.
    \noindent
    Se notará como sigue:
    \begin{equation*}
        X \leadsto U(x_1,\dots, x_n)
    \end{equation*}
\end{definicion}

\begin{prop}
    Sea $X$ una variable aleatoria tal que $X\leadsto U(x_1,\dots, x_n)$. Entonces, su función masa de probabilidad es:
    \begin{equation*}
        f(x) = P[X=x] = \left\{\begin{array}{cc}
            \nicefrac{1}{n} & x=x_i, i=1,\dots, n \\
            0 & \text{en caso contrario}
        \end{array}\right.
    \end{equation*}
\end{prop}

\noindent
Cuando $Re_X = \{1, 2, \ldots, n\}$, veamos algunos casos particulares de series:
\begin{equation*}
    \begin{split}
        E[X] &= \dfrac{1}{n}\sum_{i=1}^n i = \dfrac{1}{n}\dfrac{n(n+1)}{2} = \dfrac{n+1}{2}\\
        E[X^2] &= \dfrac{1}{n}\sum_{i=1}^n i^2 = \dfrac{1}{n}\dfrac{n(n+1)(2n+1)}{6} = \dfrac{(n+1)(2n+1)}{6}\\
        Var(X) &= \dfrac{(n+1)(2n+1)}{6} - \dfrac{(n+1)^2}{4} = \dfrac{n^2-1}{12}
    \end{split}
\end{equation*}


\subsection{Distribución binomial}
\begin{definicion}[Distribución binomial]
    Se dice que una variable aleatoria $X$ sigue una distribución binomial de parámetros $n$ y $p$, $n\in \mathbb{N}$, $p\in ]0,1[$ si modela el número de éxitos en $n$ repeticiones independientes de un ensayo de Bernouilli con probabilidad $p$ de éxito, manteniéndose esta constante en las $n$ repeticiones. Se notará como $$X\leadsto B(n,p)$$
\end{definicion}

\begin{observacion}
    Es fácil ver que la distribución de Bernouilli es un caso particular de una distribución binomial, considerando $n=1$.
\end{observacion}

\begin{prop}
    Sea $X\leadsto B(n,p)$. Entonces, la \tbf{función masa de probabilidad} de la distribución binomial es:
    \begin{equation*}
        f(x) = \binom{n}{x}p^x (1-p)^{n-x} \qquad x\in \{0,\dots,n\}
    \end{equation*}
\end{prop}

\begin{prop}
    Sea $X\leadsto B(n,p)$. Entonces, tenemos que:
    \begin{equation*}
        E[X] = np
    \end{equation*}
\end{prop}

\begin{prop}
    Sea $X\leadsto B(n,p)$. Entonces, tenemos que:
    \begin{equation*}
        \Var[X] = np(1-p)
    \end{equation*}
\end{prop}

\subsection{Distribución binomial negativa}
\begin{definicion}[Distribución binomial negativa]
    Se dice que una variable aleatoria $X$ sigue una distribución binomial negativa de parámetros $r$ y $p$, $r\in \mathbb{N}$, $p\in ]0,1[$ si modela el número de fracasos antes de llegar al $r-$ésimo éxito en un ensayo de Bernouilli con probabilidad $p$ de éxito, manteniéndose esta constante en todas las repeticiones. Se notará como $$X\leadsto BN(r,p)$$
\end{definicion}
\begin{observacion}
    Es fácil ver que la Distribución Geométrica es un caso particular de una distribución binomial negativa, considerando $r=1$.
\end{observacion}

Razonemos la función masa de probabilidad. Como la variable $X$ modela el número de fracasos antes de llegar al $r-$ésimo éxito, tenemos que en total se producen $x$ fracasos y $r$ fracasos. Como la probabilidad se mantiene constante, tenemos que la probabilidad para cierta ordenación de los sucesos es $p^r(1-p)^{x}$. No obstante, las distintas formas de reorganizar los sucesos son combinaciones de $x$ fracasos y $r-1$ éxitos\footnote{El éxito $r-$ésimo no se añade porque su posición está fijada, ya que debe ser el último suceso.}; es decir, combinaciones sin repetición. Por tanto, tenemos que en total hay $\binom{x+r-1}{x}$ formas distintas. Por tanto, a priori deducimos que la función masa de probabilidad de la distribución binomial negativa es:
\begin{equation*}
    f(x) = \binom{x+r-1}{x}(1-p)^{x}p^r \qquad x\in \mathbb{N}\cup \{0\}
\end{equation*}

No obstante, hay una expresión alternativa de la función masa de probabilidad muy útil para las demostraciones. Para ello, se introduce la siguiente definición:


\begin{definicion}
    Dado $r\in \mathbb{R}, x\in \mathbb{N}$, se define\footnote{El número combinatorio está ya definido, pero con más restricciones. Esta definición concuerda con la anterior en los supuestos anteriores.} el siguiente número combinatorio:
    \begin{equation*}
        \binom{-r}{x} = \frac{(-r)(-r-1)\cdots (-r-x+1)}{x!};~\binom{-\alpha}{0}=1, \qquad \forall r\in \mathbb{R}, x\in \mathbb{N}
    \end{equation*}
\end{definicion}


\begin{prop}
    Sea $X\leadsto BN(r,p)$. Entonces, la \tbf{función masa de probabilidad} de la distribución binomial negativa es:
    \begin{equation*}
        f(x)
        = \binom{x+r-1}{x}(1-p)^{x}p^r
        = \binom{-r}{x}(p-1)^xp^r
        \qquad x\in \mathbb{N}\cup \{0\}
    \end{equation*}
\end{prop}


\begin{prop}
    Sea $X\leadsto BN(r,p)$. Entonces, tenemos que:
    \begin{equation*}
        E[X]=\frac{r(1-p)}{p}
    \end{equation*}
\end{prop}

\begin{lema}
    Sea $X\leadsto BN(r,p)$. Entonces, tenemos que:
    \begin{equation*}
        E[X^2]=\frac{r(r+1)(1-p)^2}{p^2} + \frac{r(1-p)}{p}
    \end{equation*}
\end{lema}

\begin{coro}
    Sea $X\leadsto BN(r,p)$. Entonces, tenemos que:
    \begin{equation*}
        \Var[X]=\frac{r(1-p)}{p^2}
    \end{equation*}
\end{coro}

\subsection{Distribución de Poisson}
\begin{definicion}[Distribución de Poisson]
    Sea $X$ una variable aleatoria. Se dice que $X$ sigue una distribución de Poisson si representa el número de de ocurrencias de un determinado suceso durante un periodo de tiempo fijo o una región fija del espacio.

    Ha de cumplir las siguientes condiciones:
    \begin{enumerate}
        \item El número de ocurrencias en un intervalo o región específicas ha de ser independiente de las ocurrencias en otras zonas o intervalos de tiempo.

        \item Si se considera un intervalo de tiempo muy pequeño (o una región muy pequeña):
        \begin{itemize}
            \item La probabilidad de una ocurrencia es proporcional a la longitud del intervalo (el volumen de la región).
            \item La probabilidad de dos o más ocurrencias es despreciable.
        \end{itemize}
    \end{enumerate}

    Dicho de otra manera, si de media se tiene que en determinado intervalo se producen $\lambda$ ocurrencias, tenemos que sigue una distribución de Poisson de parámetro $\lambda$, notado por:
    \begin{equation*}
        X\leadsto \mathcal{P}(\lambda)
    \end{equation*}
\end{definicion}

\begin{prop} Sea $X\leadsto \mathcal{P}(\lambda)$. Entonces, su \tbf{función masa de probabilidad} es:
    \begin{equation*}
        f(x)=e^{-\lambda}\frac{\lambda^x}{x!}
    \end{equation*}
\end{prop}

\begin{prop}
    Sea $X\leadsto \mathcal{P}(\lambda)$. Entonces, tenemos que:
    \begin{equation*}
        E[X]= \lambda
    \end{equation*}
\end{prop}

\begin{lema}
    Sea $X\leadsto \mathcal{P}(\lambda)$. Entonces, tenemos que:
    \begin{equation*}
        E[X^2]= \lambda^2+\lambda
    \end{equation*}
\end{lema}

\begin{coro}
    Sea $X\leadsto \mathcal{P}(\lambda)$. Entonces, tenemos que:
    \begin{equation*}
        \Var[X]= \lambda
    \end{equation*}
\end{coro}

\section{Distribuciones Continuas}
\noindent
En esta sección, recordaremos algunas de las distribuciones de probabilidad continuas más importantes que se utilizan en estadística y probabilidad.


\subsection{Distribución Uniforme Continua}

\begin{definicion}[Distribución Uniforme Continua]
    Una variable aleatoria continua $X$ sigue una distribución uniforme en el intervalo $[a,b]$, con $a,b\in \bb{R}$, $a<b$, si su función de densidad
    toma un valor constante en dicho intervalo, siendo nula fuera de él.
    Lo denotaremos como $X\sim \cc{U}(a,b)$.
\end{definicion}

\begin{prop}
    Sea $X\sim \cc{U}(a,b)$. Entonces, se tiene que:
    \Func{f}{\bb{R}}{[0,1]}{x}{\begin{cases}
        \dfrac{1}{b-a} & x\in [a,b]\\
        0 & x\notin [a,b]
    \end{cases}}
\end{prop}

\begin{prop}
    Sea $X\sim \cc{U}(a,b)$, entonces su función de distribución es:
    \Func{F_X}{\bb{R}}{[0,1]}{x}{\begin{cases}
        0 & x < a\\
        \dfrac{x-a}{b-a} & x\in [a,b]\\
        1 & x > b
    \end{cases}}
\end{prop}

\noindent
Calculemos ahora los momentos de una variable aleatoria $X$ tal que $X\sim \cc{U}(a,b)$.
\begin{prop}
    Sea $X$ una variable aleatoria tal que $X\sim \cc{U}(a,b)$. Entonces, su momento no centrado de orden $k\in \bb{N}\cup \{0\}$ es:
    \begin{equation*}
        m_k = E[X^k] = \dfrac{b^{k+1} - a^{k+1}}{(k+1)(b-a)}
    \end{equation*}
\end{prop}

\noindent

Como consecuencia, tenemos que:
\begin{equation*}
    m_1 = E[X] = \dfrac{b^2 - a^2}{2(b-a)} = \dfrac{b+a}{2}
\end{equation*}


\subsection{Distribución Normal}

\noindent
Esta es la distribución de probabilidad más importante en la Teoría de la Probabilidad y la Estadística Matemática.

\begin{definicion}[Distribución Normal]
    Una variable aleatoria continua $X$ sigue una distribución normal con parámetros $\mu\in \bb{R}$ y $\sigma^2 \in \bb{R}^+$, si su función de densidad es:
    \begin{equation*}
        f_X(x) = \dfrac{1}{\sqrt{2\pi}\sigma} e^{-\dfrac{(x-\mu)^2}{2\sigma^2}}, \quad x\in \bb{R}
    \end{equation*}
    donde $\sigma=+\sqrt{\sigma^2}$.
    Lo denotaremos como $X\sim \cc{N}(\mu,\sigma^2)=\mathcal{N}_{\mu, \sigma^2}$.
\end{definicion}

A pesar de darse como definición, hemos de demostrar que efectivamente es una función de densidad.
Para ello, incluimos el siguiente Lema, cuya demostración no se incluye por su complejidad,
al requerir de integración en varias variables con cambio a coordenadas polares.
\begin{lema}[Integral de Gauss]
    Sea $a\in\bb{R}^+$, $b\in \bb{R}$. Entonces:
    \begin{equation*}
        \int_{-\infty}^{\infty} e^{-a(x+b)^2} \, dx = \sqrt{\dfrac{\pi}{a}}
    \end{equation*}
\end{lema}

\subsection{Distribución Exponencial}

\begin{definicion}[Distribución Exponencial]
    Una variable aleatoria continua $X$ sigue una distribución exponencial con parámetro $\lambda\in \bb{R}^+$, si su función de densidad es:
    \begin{equation*}
        f_X(x) = \begin{cases}
            \lambda e^{-\lambda x} & x\geq 0\\
            0 & x<0
        \end{cases}
    \end{equation*}
    Lo denotaremos como $X\sim \exp(\lambda)$.
\end{definicion}
\noindent
Veamos algunas aplicaciones de esta distribución:
\begin{itemize}
    \item La distribución exponencial se utiliza para modelar el tiempo aleatorio entre dos fallos consecutivos en \emph{fiabilidad}
    o entre dos llegadas consecutivas en \emph{teoría de colas}.

    \item También se aplica en la modelización de tiempos aleatorios de supervivencia (\emph{Análisis de Supervivencia}).
    
    \item En general, $X$ suele representar un tiempo aleatorio transcurrido entre dos
    sucesos, que se producen de forma aleatoria y consecutiva en el tiempo.
    Dichos sucesos se contabilizan mediante un proceso de Poisson homogéneo.
    El parámetro $\lm$ representa la razón de ocurrencia de dichos sucesos, que
    en este caso es constante.
\end{itemize}

\begin{prop}
    Sea $X$ una variable aleatoria tal que $X\sim \exp(\lambda)$. Entonces, su función de distribución es:
    \begin{equation*}
        F_X(x) = \begin{cases}
            1-e^{-\lambda x} & x\geq 0\\
            0 & x<0
        \end{cases}
    \end{equation*}
\end{prop}

\begin{observacion}
    Dada una v.a. $X\sim \exp(\lambda)$, se tiene que 
    \begin{equation*}
        F_X(t)=P(X<t)=\int_0^t \lambda e^{-\lambda x} dx = 1-e^{-\lambda t}, \qquad t\geq 0
    \end{equation*}
\end{observacion}

\begin{prop}
    Sea $X$ una variable aleatoria de forma que $X\sim\exp(\lambda)$. Entonces, sus momentos no centrados de orden $k\in \bb{N}$ son:
    \begin{equation*}
        E[X^k] = \dfrac{k!}{\lambda^k}
    \end{equation*}
\end{prop}

\noindent
Como consecuencia, tenemos que:
\begin{equation*}
    E[X] = \dfrac{1}{\lambda} \qquad \Var[X] = E[X^2] - E[X]^2 = \dfrac{2}{\lambda^2} - \dfrac{1}{\lambda^2} = \dfrac{1}{\lambda^2}
\end{equation*}

\section{Distribuciones Reproductivas}

\begin{definicion}[Reproductividad de distribuciones]
    Una distribución de variables aleatorias se dice reproductiva si, dadas $X_1$, $X_2$, \ldots, $X_n$ variables aleatorias que sigan una distribución de dicho tipo (aunque la distribución de cada variable tenga distintos parámetros), entonces la variable aleatoria dada por:
    \begin{equation*}
        X = \sum_{i=1}^{n} X_i
    \end{equation*}
    sigue una distribución del mismo tipo.
\end{definicion}

\begin{prop}[Distribución Reproductiva - Binomial]
    Sean $X_1, \dots, X_n$ variables aleatorias independientes y $X_i\sim B(k_i, p)$ para $i=1,\dots,n$. Entonces:
    \begin{equation*}
        \sum_{i=1}^{n}X_i \sim B\left(\sum_{i=1}^{n}k_i, p\right)
    \end{equation*}
\end{prop}

\begin{prop}[Distribución Reproductiva - Poisson]
    Sean $X_1, \dots, X_n$ variables aleatorias independientes y $X_i\sim \cc{P}(\lambda_i)$ para $i=1,\dots,n$. Entonces:
    \begin{equation*}
        \sum_{i=1}^{n}X_i \sim \cc{P}\left(\sum_{i=1}^{n}\lambda_i\right)
    \end{equation*}
\end{prop}

\begin{prop}[Distribución Reproductiva - Binomial Negativa]
    Sean $X_1, \dots, X_n$ variables aleatorias independientes y $X_i\sim BN(k_i, p)$ para $i=1,\dots,n$. Entonces:
    \begin{equation*}
        \sum_{i=1}^{n}X_i \sim BN\left(\sum_{i=1}^{n}k_i, p\right)
    \end{equation*}
\end{prop}

\begin{prop}
    Sean $X_1, \dots, X_n$ variables aleatorias independientes y $X_i\sim G(p)$ para $i=1,\dots,n$. Entonces:
    \begin{equation*}
        \sum_{i=1}^{n}X_i \sim BN(n,p)
    \end{equation*}
\end{prop}


\begin{prop}[Distribución Reproductiva - Normal]
    Sean $X_1, \dots, X_n$ variables aleatorias independientes y $X_i\sim N(\mu_i, \sigma_i^2)$ para $i=1,\dots,n$. Entonces:
    \begin{equation*}
        \sum_{i=1}^{n}X_i \sim N\left(\sum_{i=1}^{n}\mu_i, \sum_{i=1}^{n}\sigma_i^2\right)
    \end{equation*}
\end{prop}





